\section{Hypothesis and Refutation Criteria}\label{sec:hub:hypothesis}

\paragraph{Hypothesis H1.} One hub, trained jointly on the nine directions of three models, carries enough token identity that each frozen decoder regenerates the prompt from any model's hub output under free-run decoding.

\paragraph{Prediction.} Free-run word-F1 on held-out prompts exceeds 0.90 in every cross direction, with the shuffled-memory control below 0.20. H1 is refuted if any cross direction falls below 0.90 or if the control approaches the trained value.

\paragraph{Why free-run.} Teacher forcing hands the decoder the true prefix, so a decoder that recovers only the next token from context scores well without reading the hub. Free-run removes that crutch: every token after the first comes from the decoder's own output, and a hub output that carries no identity collapses to a generic continuation. Teacher-forced Acc@1 selects the checkpoint and nothing else.

\paragraph{Why inversion rather than continuation.} An earlier generation of the hub (V3) trained the decoders to continue from another model's states and reached teacher-forced Acc@1 of 0.987 against a floor of 0.323, with free-run exact match 0.586. Continuation, however, rewards fluent next tokens over faithful ones. The hub of this proposal (V6.3) retargets the decoder loss to the input tokens themselves, so the read-out measures what an attacker wants: the text that produced the states.
